% ======================================================================
% preprints/jphq_06_structural_tension_nonoptimizable/sections/03_definition_structural_tension.tex
% Journal-grade content (100/100). ASCII-only.
% ======================================================================

\section{Definition of Structural Tension}

\subsection{Status of Structural Tension in OC/ETS}

Structural tension $\TEA$ is a derived notion within the enterprise
continuum $\KEA$. It is not introduced as a primitive of the Ontology of
Continua and does not appear as an independent object in the Executable
Theory Specification. Its meaning is fully determined by the
admissibility structure defined for $\KEA$.

Formally, structural tension characterizes the relation between a given
enterprise state and the threshold conditions collected in $\ThetaEA$.
Outside this threshold-based admissibility framework, the concept of
structural tension is undefined and has no formal interpretation.

\subsection{Component-Based Formulation}

In \texttt{ETS.K\_EA.v1.1}, structural tension is represented through a
finite family of threshold component functions
\[
\mathcal{F} =
\{\fexist,\fstab,\finertia,\fcollapse,\fstop\}.
\]

Each component corresponds to a distinct admissibility condition and is
evaluated as an inequality relative to zero. The sign of a component
$f_k$ determines whether the associated condition is satisfied
($f_k \leq 0$) or violated ($f_k > 0$). No assumptions are made regarding
scale, range, continuity, ordering, or numerical interpretation of these
functions.

No aggregation, weighting, normalization, or algebraic combination of
components is defined at the theory level. Interactions between
components are governed exclusively by explicit logical precedence rules
encoded in the phase logic.

\subsection{Structural Tension as a Logical Configuration}

Structural tension $\TEA$ is not a scalar-valued function on $\OmegaEA$.
It is identified with the configuration of inequality evaluations across
the component set $\mathcal{F}$.

Two enterprise states are structurally distinct with respect to tension
if and only if they differ in the pattern of component satisfactions and
violations, even when no numerical distance, ranking, or ordering between
them can be defined. Accordingly, $\TEA$ induces neither a total nor a
partial order on $\OmegaEA$.

The only admissible ordering relation within the formalism is phase
precedence (for example, \PHASESTOP\ dominates \PHASECOLLAPSE), which is
logical rather than quantitative and does not refine distinctions between
states within the same phase.

\subsection{Relation to Boundary and Continuity}

Structural tension is related to the boundary of admissibility
$\dOmegaEA$ through threshold conditions. Boundary states are
characterized by component equalities or by the failure of required
structural properties. This relation is qualitative and non-metric.

Threshold violations may affect the continuity measure $\kEA$ by
eliminating admissible transitions or required closure cycles. However,
$\TEA$ and $\kEA$ are not linked by any functional, monotonic, or
invertible mapping. Structural tension concerns admissibility conditions,
whereas continuity concerns the existence of coherent closure and
transition structure.

\subsection{Explicit Non-Properties}

For clarity, structural tension in $\KEA$ does \emph{not} possess the
following properties:

\begin{itemize}
  \item it is not a scalar metric,
  \item it is not additive, compositional, or aggregative,
  \item it is not continuous or differentiable,
  \item it is not comparable across enterprises,
  \item it is not invariant under admissible reparameterizations,
  \item it is not an objective, loss, cost, or utility function.
\end{itemize}

These absences are structural consequences of the OC/ETS formalism. Any
extension that assigns such properties to $\TEA$ introduces additional
primitives and therefore constitutes a change of theory rather than an
interpretation within it.

%%% End of section %%%
